% Solutions to Chapter 4, from lectures/L04/appendix/solutions: the README's prose, and the two
% reference programs typeset straight from their files.

\section{Chapter 4: The Factory Pattern}
\label{sol:sec:c4}
The exercises in \secref{c4:sec:exercises} build one program twice: once with a factory that
returns raw pointers (Sets~1 to~3), and once with a factory that returns \code{std::unique_ptr}
(Set~4). The reference solutions are those two programs, in
\file{lectures/L04/appendix/solutions/raw_factory} and
\file{lectures/L04/appendix/solutions/smart_factory}. Each directory has a Makefile, so
\code{make} there builds the program and runs it; the application logic runs forever, so stop it
with Ctrl+C.

\note Try to solve the exercises yourself before looking at the solutions.

\subsection{The raw-pointer factory}
\label{sol:set:c4.1}\label{sol:set:c4.2}\label{sol:set:c4.3}
Application logic implemented using raw pointers and a raw-pointer factory: the solution to
Exercises~\ref{c4:ex:1.1} to~\ref{c4:ex:3.2}.

\paragraph{Serial interface and drivers}
Exercises~\ref{c4:ex:1.1} to~\ref{c4:ex:1.3}.

\cppfile{lectures/L04/appendix/solutions/raw_factory/include/driver/serial/interface.hpp}
\cppfile{lectures/L04/appendix/solutions/raw_factory/include/driver/serial/stub.hpp}
\cppfile{lectures/L04/appendix/solutions/raw_factory/include/driver/serial/esp32s3.hpp}

\paragraph{Factory interface with raw pointers}
Exercises~\ref{c4:ex:2.1} to~\ref{c4:ex:2.3}.

\cppfile{lectures/L04/appendix/solutions/raw_factory/include/driver/factory/interface.hpp}
\cppfile{lectures/L04/appendix/solutions/raw_factory/include/driver/factory/esp32s3.hpp}
\cppfile{lectures/L04/appendix/solutions/raw_factory/include/driver/factory/stub.hpp}

\paragraph{Application logic with raw pointers}
Exercises~\ref{c4:ex:3.1} and~\ref{c4:ex:3.2}.

\cppfile{lectures/L04/appendix/solutions/raw_factory/include/app/logic/logic.hpp}
\cppfile{lectures/L04/appendix/solutions/raw_factory/source/main.cpp}

\subsection{The smart-pointer factory}
\label{sol:set:c4.4}
Application logic implemented using smart pointers and a smart-pointer factory: the solution to
Exercises~\ref{c4:ex:4.1} to~\ref{c4:ex:4.3}.

Only the factories and the logic class change. The three serial driver headers,
\file{driver/serial/interface.hpp}, \file{driver/serial/stub.hpp} and
\file{driver/serial/esp32s3.hpp}, are identical to the raw-pointer versions above, and
\file{main.cpp} differs from its raw-pointer version only in its file comment. The four files that
do change are:

\cppfile{lectures/L04/appendix/solutions/smart_factory/include/driver/factory/interface.hpp}
\cppfile{lectures/L04/appendix/solutions/smart_factory/include/driver/factory/esp32s3.hpp}
\cppfile{lectures/L04/appendix/solutions/smart_factory/include/driver/factory/stub.hpp}
\cppfile{lectures/L04/appendix/solutions/smart_factory/include/app/logic/logic.hpp}

\subsection{Selecting the factory}
\label{sol:c4:selecting}
The reference solution uses the \code{STUB} macro in \file{main.cpp} to switch between the ESP32-S3
factory and the stub factory. This is only for convenience and was not required by the exercise.

The example application can be built using either the ESP32-S3 drivers or the stub drivers.

By default, the stub drivers are used due to the macro \code{STUB} being defined in
\file{main.cpp}:

\begin{cppcode}
#define STUB
\end{cppcode}

\noindent To use the ESP32-S3 drivers instead, comment out or remove the definition:

\begin{cppcode}
// #define STUB
\end{cppcode}

\noindent The application logic (\code{app::logic::Logic}) does not need to be modified. Only the
selected factory changes.

\subsection{Reflection}
\label{sol:c4:reflection}
The answers to the reflection questions in Exercise~\ref{c4:ex:3.2}.

\textbf{What code in \code{Logic} had to change when switching from the ESP32-S3 factory to the stub
factory?}
\begin{itemize}
  \item No code in \code{Logic} had to change.
  \item The logic class depends only on the abstract factory interface and the abstract serial
    interface, not on the concrete ESP32-S3 or stub implementations.
  \item This means the application logic can use either real hardware drivers or stub drivers
    without being modified.
\end{itemize}

\textbf{What code in \code{main()} had to change?}
\begin{itemize}
  \item Only the concrete factory type created in \code{main()} had to change.
  \item Instead of creating a \code{driver::factory::Esp32s3}, \code{main()} creates a
    \code{driver::factory::Stub}.
  \item The rest of the code can stay the same, because both factories implement the same factory
    interface.
\end{itemize}
